% 本文件由 LaTeX 排版 Agent 根据有效 translation.json 自动生成。
% author=Jerome; work_id=3003; title=圣希拉里翁传

\chapter{\WorkTitle{圣希拉里翁传}}

% structure: p0001 source=h1 target=omitted reason=page-title-already-rendered-by-parent

圣希拉里翁传由\Person{热罗尼莫}于390年在\Place{白冷}撰写。其目的在于提倡他所献身的苦修生活。书中虽包含许多传说，但也有若干与真实历史相联系的记述，无论如何，它是对四世纪人类心智状态的一个奇异记录。\Place{德国}兴起的一种理论认为它是一部宗教传奇，似乎毫无根据。可能，按\Person{热罗尼莫}的意图，它是为教会历史所作的一项贡献，他本拟撰写教会历史但未实现。（见\WorkTitle{圣玛尔基传}第一章。）\par

1. 在我开始撰写真福\Person{希拉里翁}的生平之前，我恳求曾居住在他内的\OriginalTerm{圣神}{Holy Spirit}的助佑，愿那赐予圣者德能的那一位，赐给我言词的能力来叙述它们，使我的言语能与他的事迹相称。因为那些成就伟业之人的德能，其受推崇的程度，取决于天才人物赞颂它的能力。\Place{马其顿}的\Person{亚历山大大帝}，即\Person{达尼尔}所论及的公绵羊、豹或公山羊，来到\Person{阿喀琉斯}的墓前，喊道：\Quote{幸福的青年！能有伟大的传扬你价值的人。}他当然是指\Person{荷马}。然而，我不得不讲述一位如此卓著之人的生平与言行，即使\Person{荷马}在此，他要么会嫉妒我的题材，要么会证明自己力有不逮。诚然，那位圣者，\Place{塞浦路斯}\Place{萨拉米斯}的主教\Person{厄丕法尼乌斯}，曾与\Person{希拉里翁}多有往来，在写给众人的一封短函中赞扬了他，但这封信流传甚广。然而，笼统地赞美死者是一回事，叙述他们特有的德行又是另一回事。因此我们接手他开创的工作，是服事他而非亏负他：我们鄙视某些人的毁谤，他们曾诋毁我的英雄\Person{保禄}，如今或许要诋毁\Person{希拉里翁}；前者因独修生活受责备；后者则因其与世交往而受非难；前者总是隐而不见，所以他们以为他并不存在；后者为众人所见，因而被视为无足轻重。这正是他们的祖先\Person{法利塞人}古时所行的！他们不喜悦（\BibleRef{玛 11:18}）在旷野中禁食的\Person{若翰}，也不喜悦在繁忙人群中又吃又喝的主与救主。但我将着手我已决定的工作，继续前行，对\Person{斯库拉}之犬的吠叫充耳不闻。\par

2. \Person{希拉里翁}的出生地是\Place{塔巴塔}村，位于\Place{巴勒斯坦}\Place{迦萨}城以南约五英里。他的父母是崇拜偶像者，因此，正如俗语所说，玫瑰长在荆棘上。他们把他托付给\Place{亚历山大}的一位文法教师照管，在那里，就其年龄所能，他显示出非凡的才能和品性；不久便受到众人喜爱，成为出色的演说家。比这一切更重要的是，他信从主耶稣，不沉迷于竞技场的疯狂、斗兽场的血腥、剧院的放纵；他全部的喜乐在于教会的聚会。\par

3. 那时他听说了\Person{圣安当}的盛名，\Person{安当}的名字传遍了\Place{埃及}各族。他燃起见他的渴望，便出发前往旷野。他一见到\Person{安当}，就改变了自己从前的生活方式，与他一同住了大约两个月，学习他生活的方式和行为的庄重：他祈祷的恒心，与弟兄来往时的谦逊，斥责时的严厉，劝诫时的热忱。他也注意到，这位圣人从不因身体的软弱而打破节制的规矩，或偏离饮食的简朴。最后，他再也无法忍受那些因各种疾病或魔鬼攻击而来访圣人的群众，认为在旷野中忍受城市的拥挤实属怪异，便想最好像\Person{安当}开始那样开始。他说：\Quote{\Person{安当}就像一个证明了自己勇气的英雄，正在收获胜利的果实。我还没有踏上士兵的征程。}于是他同几位隐修士返回家乡，那时他的父母已去世，他把部分财产给了兄弟，部分给了穷人，自己什么也不留，因为他敬畏地记起\BibleBook{宗徒}{大事录}中的那段话，并害怕\Person{阿纳尼雅}和\Person{撒斐辣}的榜样和惩罚；他尤其牢记主的话：\BibleRef{路 14:33}\BibleQuote{你们中不论是谁，如不舍弃一切所有，不能做我的门徒。}当时他大约十五岁。因此，赤裸裸地，并武装着基督的武器，他进入了旷野，该旷野从\Place{迦萨}的港口\Place{马约马}向左延伸七英里，沿着海岸通往\Place{埃及}。虽然那地方有抢劫和流血的记录，他的亲属和朋友警告他正面临的危险，但他藐视死亡，为能逃避死亡。\par

4. 他的勇气和年幼本应使众人惊奇，若不是他内心炽热，他的眼睛闪耀着信德的光芒和火花。他的面颊光洁，身体瘦弱纤细，经不起冷热所能加予的最轻微伤害。那么，怎么办呢？他除了穿一件苦衣衬衣，以及真福\Person{安当}在他出发时赠送的一件皮斗篷，和最粗糙的一条毯子外，别无遮体之物，却以一边是海、另一边是沼泽的广大可怖的旷野为乐。他的食物只是日落后十五颗干无花果。又因那地区以强盗闻名，他的习惯是从不在同一地方久留。魔鬼该怎么办？它能转向何处？它曾自夸说：\BibleRef{依 14:14}\BibleQuote{我要升到天上，我要把我的宝座置于天空的星辰之上，我要与至高者相似。}现在却看见自己被一个少年征服并践踏在脚下，这少年的年龄尚不容许犯罪。\par

5. 因此撒殫搔養他的感官，照其所常為的，在他日趨成熟的身體內點燃情慾之火。這位基督學校中的初學者，被迫思想他所不知的事，並在心中盤旋他從未經驗過的全套想法。他對自己發怒，捶打胸膛（彷彿用手的一擊就能驅逐思想），喊道：\Quote{驢！} \Quote{我要止住你的踢蹬，我不再以大麥餵你，而是以糠秕。我要以飢渴削弱你，我要以重擔壓傷你，我要驅你經過炎熱與寒冷，使你更思念食物而非淫亂。} 此後三四天，他以草汁和幾粒乾無花果維持下沉的精神，屢次祈禱、歌唱，並鋤地，使禁食的痛苦因勞苦而加倍。同時他用燈心草編織籃子，效法埃及隱修士的紀律，實踐宗徒的教訓：\BibleQuote{若有人不願意工作，就不應當吃飯。} \BibleRef{得後 3:10} 由於這些實踐，他變得如此衰弱，身體如此消瘦，以致骨骸幾乎無法相連。\par

6. 一夜，他開始聽見嬰兒的哀哭、羊群的咩叫、牛群的哞鳴、似是婦女的哀嘆、獅子的吼叫、軍隊的喧囂，此外還有各種奇異的呼喊，使他在驚恐中尚未看見就畏縮其聲。他明白這是魔鬼在戲弄，便雙膝跪下，在額上劃十字聖號。如此武裝躺臥，他更加勇敢地作戰，半渴望看見那些他聽來戰慄者，焦慮地四處張望。與此同時，在明亮的月光下，他忽然看見一輛戰車疾馳而來。他呼求耶穌，剎那間在他眼前，大地裂開，整隊人馬被吞沒。於是他說：\BibleQuote{馬和騎馬的，祂都拋在海中。} \BibleRef{出 15:1} 又說：\BibleQuote{有人靠車，有人靠馬，但我們要因我們上主天主的名而誇勝。}\par

7. 他的誘惑如此之多，魔鬼的網羅晝夜如此多變，若我願一一敘述，一卷書也不夠。多少次當他躺臥時，裸女向他顯現；多少次當他飢餓時，豐盛的筵席出現！有時他祈禱時，一匹嗥叫的狼跳過，或一頭咆哮的狐狸；他歌唱時，眼前出現角鬥表演，一個剛被殺死的人似乎倒在他腳下，求他埋葬。\par

8. 有一次，他伏地祈禱。如人之常情，他的注意力從祈禱中游移，思想別事，此時一個折磨者跳到他背上，用腳跟踢他的肋旁，用馬鞭抽打他的頸項，喊道：\Quote{來！你為何睡覺？} 接著大聲笑著問他是否累了，想要一些大麥。\par

9. 從十六歲到二十歲，他在一間用蘆葦和莎草搭成的小屋中遮擋炎熱和雨水。後來他為自己建造了一間小室，至今猶存，高五呎，比他自己的身高還矮，長度也僅多一點。人可能會以為那是一間墳墓，而非房屋。\par

10. 他每年在復活節剃髮一次，直到去世都習慣睡在光地上或用燈心草鋪的床上。他穿上的苦衣從不洗滌，並常說在山羊毛衣中追求清潔太過分了。他也不更換襯衫，除非所穿的那件幾乎破爛。他將聖經牢記在心，祈禱和歌唱後，常好像在天主面前朗誦。若要一一敘述他靈性攀登的每個步驟，將是冗長乏味的；因此我將以摘要方式向讀者陳述，描述他在各階段的生活方式，之後再回到適當的歷史順序。\par

11. 從二十歲到二十七歲，三年間他的食物是半品脫用冷水浸濕的小扁豆；其後三年是乾麵包加鹽和水。從二十七歲到三十歲，他以野草和某些灌木的生根為生。從三十一歲到三十五歲，他吃六盎司大麥麵包和少許無油煮過的蔬菜。但發現視力模糊，全身因疥瘡和乾癬而萎縮，他便在食物中加了油，直到六十三歲都遵循這適度的路線，未嘗任何水果、豆類或其他任何東西。後來見身體衰弱，以為死亡將近，從六十四歲到八十歲，他戒絕了麵包。他的精神熱忱如此驚人，以至於在別人通常對生活稍加放縱的時候，他卻像初學者一樣開始事奉主。他用麵粉和搗碎的草製成一種稀飯，食物和飲料合計不足六盎司；在遵守這飲食規則時，他從不在日落前進食，即使在節日或重病時也不例外。但現在是回到事件進程的時候了。\par

12. 當他仍住在那小屋時，十八歲那年，有盜賊夜間來找他，或是以為他有什麼可以拿走，或是認為若一個孤獨少年不怕他們的襲擊，他們就會被輕視。他們從傍晚到天亮在海和沼澤之間來回搜尋，卻找不到他的住處。然後，在日光下發現這少年，他們半開玩笑地問他：\Quote{若盜賊到你這裡來，你會怎麼做？} 他回答：\Quote{一無所有的人不怕盜賊。} 他們說：\Quote{無論如何，你可能會被殺。} \Quote{我可能，} 他說，\Quote{我可能；因此我不怕盜賊，因為我準備好去死。} 於是他們驚奇他的堅定和信德，承認他們如何在夜間迷路，眼睛如何被弄瞎，並應許將來度更嚴謹的生活。\par

13. 他在旷野已度过了二十二年，成为巴勒斯坦各城街谈巷议的话题，尽管众人只是闻名而已。第一个敢于闯入真福希拉利翁面前的人，是厄娄特鲁波利的一位妇人，她因不孕而被丈夫轻视（因她十五年来未生育）。他未曾预料她的到来，她却突然俯伏在他脚前。\Quote{请原谅我的冒昧，}她说，\Quote{请怜悯我的急需。你为何转脸不看？为何拒绝我的恳求？不要视我为女子，而当作怜悯的对象。是女人的性生育了救主。\BibleRef{路 5:31}健康的人不需要医生，有病的人才需要。}最后，过了许久，他不再转脸，而是看着那妇人，问她来的缘由和流泪的原因。得知后，他举目望天，命她要有信心，然后在她离去时为她流泪。一年之内，他见到她带着一个儿子。\par

14. 这第一个奇迹之后，又发生了另一个更大更引人注目的奇迹。阿利斯泰奈特是厄尔丕狄的妻子（厄尔丕狄后来成为执政官），她在本族人中颇有名望，在基督徒中更为知名；她与丈夫拜访了真福安当后，在返回途中，因三个孩子的疾病而滞留在迦萨。在那里，无论是由于空气污浊，还是如后来所显明的，是为了光荣天主的仆人希拉利翁，三个孩子都患上了半三日热，被医生们放弃。母亲躺着哀哭，或者可以说，在三个儿子的尸体之间来回踱步，不知该先为哪一个哀悼。当她得知附近旷野有一位隐修士时，便忘记了主母的身份（她只记得自己是一位母亲），带着侍女和太监出发，丈夫好不容易才说服她骑上一头驴。到达圣人后，她说：\Quote{我因我们最仁慈的天主耶稣求你，我因他的十字架和宝血恳求你，将我的三个儿子还给我，好使我们主救主的名在外邦人的城中受显扬。那时，他的仆人将进入迦萨，偶像玛尔纳将倒在地上。}起初他拒绝，说他从未离开过自己的小屋，也不习惯进入房屋，更不用说城市了；但她扑倒在地，反复哭喊：\Quote{希拉利翁，基督的仆人，把我的孩子还给我：安当在埃及保全了他们，你在叙利亚拯救他们吧。}所有在场的人都哭了，圣人自己也一边流泪一边拒绝了她。还需要多说吗？那妇人一直不离开，直到他答应日落之后进入迦萨。他来到那里，在每张床和发烧的四肢上划十字圣号，并呼求耶稣之名。圣名何其灵验！仿佛从三个泉源同时涌出汗水：就在那一刻，他们进食了，认出了哀哭的母亲，并感谢天主，热切地亲吻圣人的手。此事传开后，名声远播，人们从叙利亚和埃及蜂拥而至，以致许多人信了基督，并自称隐修士。因为在巴勒斯坦还没有隐修院，在叙利亚也无人认识隐修士，直到圣人希拉利翁。正是他开创了这种生活方式和虔敬，并首次在该省训练人们如此生活。主耶稣在埃及有年迈的安当，在巴勒斯坦有年轻的希拉利翁。\par

15. 法奇狄亚是埃及城市里诺科路拉的一个小村庄。从这个村庄，一位失明了十年的妇人被带到真福希拉利翁面前，由弟兄们（因为当时已有许多隐修士与他在一起）引见给他，她声称已经将全部财产花费在医生身上。圣人回答说：\Quote{你若将浪费在医生身上的给了穷人，真正的医生耶稣早就医治了你。}但她大声哭喊，恳求怜悯，他便效法救主，向她的眼睛吐唾沫，同样立即见效。\par

16. 还有一位来自迦萨的马车夫，在他的战车上被魔鬼击打，全身僵硬，既不能动他的手，也不能弯他的颈。他被放在担架上抬来，只能通过移动舌头来表示他的请求。他被告知，除非他首先信仰基督并承诺放弃他以前的职业，否则不能痊愈。他信了，承诺了，便痊愈了：他为灵魂的得救比身体的得救更加欢欣。\par

17. 一位非常强壮的年轻人，名叫玛尔西塔斯，来自耶路撒冷附近，对自己的力量极为自负，曾扛着十五蒲式耳谷物走很长距离，并认为能胜过驴子的耐力是他最高的荣耀。此人被一个凶恶的魔鬼折磨，不能忍受链子或脚镣，甚至打破门闩和门闩。他曾咬掉许多人的鼻子和耳朵，打碎一些人的脚，另一些人的腿。他使众人如此恐惧，以至于被锁上链子，从四面八方用绳子像野牛一样拖到隐修院。弟兄们一看见他，就大为惊恐（因那人身材巨大），并告诉了父亲。他坐着，命令把他带到他面前并释放。当他被释放时，\Quote{低下头，}他说，\Quote{过来。}那人开始颤抖；他扭过头，不敢正视他的脸，但放下了所有的凶猛，开始舔他坐着的脚。最后，附在年轻人身上的魔鬼被圣人的驱魔折磨，在第七天出来了。\par

18. 我们也不可略而不提\Person{奥里翁}的事。他是\Place{红海}海岸\Place{艾拉}的一位首领、富有的市民，被一群魔鬼附身，被带到圣人面前。他的手、颈、肋、脚都缠着铁链，他的怒目预示着疯狂发作。圣人正与弟兄们同行，讲解一段圣经时，那人挣脱了看守者的手，从后面抱住他，将他举起。众人大喊，因为他们害怕他可能压坏圣人因禁食而消瘦的肢体。圣人微笑着，说：\Quote{安静，让我独自与我的角力对手较量。}于是他将手从肩后弯回去，直到触到那人的头，抓住他的头发，把他拉转过来，使他与他脚对脚；然后他伸直双手，用自己的双脚踩住那人的双脚，同时连连喊道：\Quote{折磨你！你们这群魔鬼，折磨！}那受苦者大声喊叫，把头向后弯直到触地，而圣人说：\Quote{主耶稣，释放这个可怜的人，释放这个俘虏。战胜许多人，如同战胜一个人一样，是祢的权能。}我现在要讲的是无与伦比的事：从一个人的口中听到了不同的声音，仿佛是众人的喧嚷。好吧，他也被治愈了，不久后带着妻子和孩子们来到隐修院，带来许多礼物表达他的感激。圣人这样对他说：\Quote{你没有读过\Person{革哈齐}和\Person{西满}所遭遇的事吗？一个接受了报酬，另一个给出了报酬，前者是为了出卖恩宠，后者是为了购买恩宠。}\Person{奥里翁}流着泪说：\Quote{请收下，分给穷人吧。}他回答说：\Quote{你最善于分配你自己的礼物，因为你行走在城市的大街小巷，认识穷人。我为什么舍弃自己的东西，却要寻求别人的呢？对许多人来说，穷人的名字是他们贪婪的借口；但怜悯没有诡计。没有人比那什么都不为自己保留的人施舍得更好。}那人忧愁地伏在地上。他说：\Quote{不要忧愁，我儿；我为自己做的，也是为你做的。如果我收了这些礼物，我自己会得罪天主，而且那魔鬼群还会回到你身上。}\par

19. 有一个关于\Place{加沙}的马约米特人的故事，不可默然不提。他在离隐修院不远的海岸开采建筑石材时，完全瘫痪了，被他的工友抬到圣人那里后，立刻完全健康地回去工作。我应该解释一下，\Place{巴勒斯坦}和\Place{埃及}的海岸天然由软沙和砾石组成，渐渐凝固变硬成为岩石；因此虽然看起来一样，摸起来却不一样了。\par

20. 另一个故事关于\Person{伊塔利库斯}，他是同城的一位市民。他是基督徒，为竞技场豢养马匹，与\Place{加沙}的杜维姆维里的马匹竞赛，那人是偶像神\Person{玛尔纳斯}的崇拜者。这一习俗至少在罗马城市中与\Person{罗慕路斯}时代一样古老，是为纪念成功掳掠萨宾妇女而设立的。战车围绕竞技场跑七圈，以尊崇\Person{孔苏斯}——他作为谋略之神。胜利属于使对手的马匹疲惫不堪的那一队。现在，\Person{伊塔利库斯}的对手雇用了一个魔法师，通过某些恶魔的咒语来激发他自己的马，并阻挡对手的马。因此\Person{伊塔利库斯}来到有福的\Person{希拉良}那里，恳求他的帮助，与其说是为了损害对手，不如说是为了保护自己。这位可敬的老人觉得在这种琐事上浪费祈祷似乎荒谬。于是他微笑着，说：\Quote{你为什么不把马匹的价钱分给穷人，为了你灵魂的救恩？}他的访客回答说，他的职务是一项公共义务，他这样做与其说是出于选择，不如说是出于强迫；没有一个基督徒可以使用魔法，宁可寻求\Person{基督}的仆人的帮助，尤其是对抗\Place{加沙}的人，他们是天主的敌人，并且会因战胜\Person{基督}的教会而狂喜，更甚于因战胜他本人。因此，在在场的弟兄们的请求下，他命人将他惯常饮用的一个陶杯盛满水，交给\Person{伊塔利库斯}。后者接过杯子，将水洒在他的马厩、马匹、战车御者、战车以及赛道的栅栏上。人群异常激动，因为对手曾嘲弄地散布了将要发生的事，而\Person{伊塔利库斯}的支持者则因他们自许的胜利而兴高采烈。信号发出；一队飞向终点，另一队停滞不动：一队的战车轮子炽热，而另一队几乎只瞥见对手飞驰而过的背影。人群的欢呼声高涨成雷鸣，连异教徒自己也一致宣称\Person{玛尔纳斯}被\Person{基督}征服了。之后，愤怒的对手要求将\Person{希拉良}作为基督徒魔法师拖去处决。这决定性的胜利以及随后在连续竞技比赛中取得的多次胜利，使许多人归向了信仰。\par

21. 在同一\Place{迦萨}集镇附近，有一个年轻人疯狂地爱上了一位天主贞女。他一再尝试那些通常导致贞洁毁灭的触碰、戏谑、点头和耳语，但通过这些方式毫无进展，便前往\Place{孟斐斯}找了一名巫师，向他诉说自己可怜的状况，然后借助巫术，返回继续向该贞女进攻。于是，在\Person{阿斯克勒庇俄斯}的司祭（这人不是医治灵魂，而是毁灭灵魂）教导一年后，他充满了先前已容许进入心灵的情欲，在女孩家的门槛下埋藏了一些写在\Place{塞浦路斯}铜板上的魔法咒语和令人憎恶的图案。随即，这少女开始显出疯狂的迹象，扔掉头巾，撕扯头发，咬牙切齿，并大声呼喊那年轻人的名字。她深沉的爱变成了一种狂乱。因此，她的父母将她带到隐修院，交给那老年圣人。刚一这样做，魔鬼就开始嚎叫并承认：\Quote{我是被迫的，是违背我的意志被拖来的。当我常常在梦中欺骗\Place{孟斐斯}人时，我是多么快活！我现在遭受何等十字架的折磨，何等的痛苦！你强迫我出去，而我却被束缚在门槛下。我出不去，除非那囚禁我的年轻人放我走。}老人回答说：\Quote{你的力量一定很大，如果一条线和一块铜板就能把你束缚住。告诉我，你怎么敢进入这属于天主的贞女？}\Quote{为的是保存她的贞洁，}他说。\Quote{你这奸污贞洁的人，还保存她！你为什么不去进入那派你来的人？}\Quote{我为什么要进入一个与我的同伴——情魔——结盟的人？}他回答道。但圣人没有下令寻找那年轻人或那符咒，直到少女经历了一个净化的过程，因为害怕会被人认为是靠咒语释放了魔鬼，或者他自己相信了魔鬼所说的话。他宣称魔鬼是狡猾的、擅长伪装的，并在贞女恢复健康后严厉斥责她，因为她的行为给了魔鬼进入的机会。\par

22. 他不仅在\Place{巴勒斯坦}及邻近的\Place{埃及}或\Place{叙利亚}城市中享有盛名，而且他的名声已传到了遥远的省份。\Person{君士坦提乌斯}皇帝的一位军官，因其金发和俊美仪表透露出他的国籍（他的国土位于撒克逊人和阿勒曼尼人之间，面积不大但强盛，史学家称之为日耳曼，但现在称为法兰克），长期以来——即从幼年起——就被一个魔鬼纠缠，魔鬼强迫他在夜间嚎叫、呻吟和咬牙切齿。于是他秘密向皇帝请求一张驿站通行证，坦白告诉他要它的原因，又获得了致\Place{巴勒斯坦}总督的书信，便以其堂皇排场和大量随从来到\Place{迦萨}。他向当地议员询问隐修士希拉流住在哪里，\Place{迦萨}人大为惊恐，以为他是皇帝派来的，便将他领到隐修院，以便向这样一位显赫人物表示敬意，同时，如果他们先前因对希拉流作恶而犯有罪过，希望借当下的恭顺来抹去它。那时老人正在软沙上散步，低声吟唱\BibleBook{圣}{咏}中的某段。看到那么大一群人走近，他停下来，向众人一一回礼，同时举手祝福，过了一小时后他让其余人退下，但让他的访客连同仆人和军官留下：因为从这人的眼神和面容，他已知晓他来的原因。天主的仆人一问，那人立即踮起脚尖，双脚几乎不触地，用叙利亚语（圣人正是用这种语言问他）以狂野的咆哮回答。纯正的叙利亚语从只懂法兰克语和拉丁语的蛮族的嘴唇中流出，没有一处咝音、送气音或巴勒斯坦语习语的缺失。魔鬼随后承认它是如何进入他体内的。此外，为了让只懂希腊语和拉丁语的口译员能听懂，希拉流又用希腊语问他，当魔鬼用同样的话语作了同样的回答，并以许多场合他曾被施以法术为由辩解，声称他不得不屈服于魔法时，圣人说：\Quote{我不在乎你是如何进来的，}圣人说，\Quote{但我以我们主耶稣基督之名命令你出来。}那人痊愈后，以粗朴的率直向他献上十镑金子。但是圣人只从他那里接受了面包，并告诉他，以这种食物为生的人视金子如泥土。\par

23. 谈论人类尚不足够；每天也有疯狂的野兽被带到他面前，其中有一只巨大的巴克特里亚骆驼，由三十多人用粗绳紧紧拉住。它已经伤害了许多人。它双眼充血，口吐白沫，肿胀的舌头翻滚，最令人恐惧的是它响亮而可怕的吼叫声。老人命令放开它。立刻，那些带它来的人和圣人的侍从都毫无例外地逃走了。圣人独自上前迎向它，用叙利亚语对它说：\Quote{你吓不到我，魔鬼，虽然你现在的身体巨大。无论是在狐狸还是骆驼中，你都是一样的。}同时，他伸出手站着。那只野兽狂暴着，看起来仿佛要吞吃依拉良，走上前来，但立刻倒下了，把头放在地上，令所有在场的人惊讶的是，它突然表现出与之前凶暴同样的温顺。但老人向他们说明，魔鬼为了人的缘故，甚至夺取驮兽；它被对人类的如此强烈仇恨所燃烧，以至于它想要毁灭的不仅是人，还有属于人的一切。作为例证，他补充说，在它被允许试探圣约伯之前，它已经毁灭了他一切所有的。也不应因主命令魔鬼杀死了两千头猪而惊惶，因为目睹奇迹的人若没有同样多的猪冲入毁灭，就无法相信那男人身上出来了如此众多的魔鬼，这表明是一支军团在驱使他们。\par

24. 若我想述说他所行的一切奇迹，时间就不够了。因为主将他提升到如此荣耀的地步，以至于连真福安当听闻他的生活后，也写信给他，并欣然收到他的回信。如果有时从叙利亚来求见安当的病人，他会对他们说：\Quote{你们何必不辞辛劳远道而来，你们那里有我的儿子依拉良呢？}然而，效法他的榜样，整个巴勒斯坦涌现出无数的隐修院，所有隐修士都涌向他。看到这情景，他赞美主的恩宠，并个别劝勉他们为灵魂的益处，告诉他们这世界的格局正在逝去，真正的生命是由现今的苦难所换得的。\par

25. 为了给隐修士们树立谦逊和热忱的榜样，他惯常在固定的日子、在葡萄收获期之前探访他们的居所。弟兄们一得知此事，便都一同来迎接他，并与他们杰出的领袖一道巡行各隐修院，带着食物，因为有时聚集了两千人之多。但随着时间推移，周围所有的居民都乐意给邻近的隐修士提供食物，以款待这些圣人。此外，他所操心的，是任何一位弟兄，无论多么卑微贫穷，都不得被忽略；这可由他曾经一次进入卡德士旷野探访他的一位门徒一事为证。他带着一大群隐修士来到埃卢萨，恰逢那日的年度庆典将所有民众聚集到维纳斯神庙。这位女神因晨星而受崇拜，萨拉森民族便奉献于此。这座城由于其地理位置，也在很大程度上是半野蛮的。因此，当听闻圣依拉良经过时（他曾多次治愈许多被魔鬼附身的萨拉森人），他们便带着妻儿成群结队地来迎接他，低着头，用叙利亚语呼喊\OriginalTerm{祝福}{Barech}，即\Emph{祝福}。他谦恭而和蔼地接待了他们，并祈祷他们崇拜天主而不是石头；同时，他泪流满面，仰望天堂，许诺如果他们相信基督，他会经常拜访他们。因着天主奇妙的恩宠，他们不让他离开，直到他画出了教堂的轮廓，并且他们的戴着花冠的司铎已被基督的标记所标记。\par

26. 又一年，当他再次出发探访隐修院，并列出他必须留宿和途经时需看望的人的名单时，隐修士们知道他们中有一位吝啬的弟兄，同时渴望治愈他的毛病，便请圣人留宿在那位弟兄处。他回答说：\Quote{你们愿意我对你们造成伤害，对那位弟兄造成烦扰吗？}那位吝啬的弟兄听到这话后感到羞愧，并在全体弟兄的大力支持下，终于从圣人那里勉强获得承诺，将他的隐修院列入歇息之所的清单。十天后，他们来到他那里，发现看守人已经守候在他们必经的葡萄园中，用石头、土块和投石器驱赶所有来者。早晨，他们都没有吃一颗葡萄就离开了，而老人微笑着，假装不知道所发生的事。\par

27. 有一次，他们被另一位名叫撒布斯的弟兄款待（当然我们不可说出那吝啬者的名字，但可以提这位慷慨者的名字），因为那天是主日，他们都受邀进入葡萄园，以便在用餐时间之前用葡萄缓解旅途的辛劳。圣人说：\Quote{谁若在灵魂的舒畅之前寻求身体的舒畅，愿他受诅咒。让我们祈祷，让我们歌唱，让我们尽我们对天主的本分，然后我们才赶到葡萄园去。}礼仪结束后，他站在一处高地，祝福了葡萄园，并让他自己的羊群去放牧。当时享用者不下三千人。而整座葡萄园原估计可产一百壶酒，但三十天内，他使它值三百壶。那位吝啬的弟兄收获比往常少得多，并且发现连他已有的也都变成了醋，他十分悲伤。老人在此事发生之前已向许多弟兄预言过了。他尤其厌恶那些因缺乏信德而为将来囤积、并操心花费、衣裳或其它那些随世界而逝去之物的隐修士。\par

28. 最后，他甚至不愿看一位住在约五里之外的弟兄一眼，因为他查明此人极度吝啬地守护自己的一小块地，还存有一点钱。这位犯错者想与老人和好，常到弟兄们那里去，尤其去找倍受希拉里翁喜爱的\Person{赫西基乌斯}。有一天，他带来一捆刚摘下的青豌豆。赫西基乌斯在傍晚将它放在桌上，老人随即喊道，他受不了那臭味，问那豌豆是从哪里来的。赫西基乌斯回答说，某位弟兄把他地里初熟的果实送给了弟兄们。\Quote{你难道没闻到，}他说，\Quote{那可怕的臭味，在豌豆里嗅到贪婪的恶臭？把它扔给牲口，扔给畜类，看看它们能不能吃。}遵照他的命令，豌豆刚被放入食槽，牲口就极度惊慌，大声吼叫，挣断缰绳，四散奔逃。原来，老人藉着恩宠，能从身体、衣服以及任何人触摸过的东西的气味，分辨出那人是受哪个魔鬼或哪种恶习的困扰。\par

29. 到了第六十三年，这位老人主持着一个大隐修院，住着许多弟兄。不断有人带着各种患病或被不洁之神附身的人前来，多到整个旷野周边都挤满了各种各样的人。圣人见此情景，天天哭泣，无比遗憾地追忆他从前的生活方式。一位弟兄问他为何如此沮丧，他回答说：\Quote{我又回到了世俗，生前就已得到了我的报酬。巴勒斯坦及邻近省份的人以为我有些了不起，借口为了管理好弟兄们而设隐修院，我却拥有了卑俗生活所需的一切。}然而，弟兄们看守着他，尤其是赫西基乌斯，他对老人怀有极其虔敬的热爱与崇敬。他在这样的哀叹中度过了两年之后，我们之前提到的那位\Person{阿里斯特内特}夫人——当时她是总督的妻子，却毫无总督夫人的浮华——前来拜访他，也打算去探望\Person{圣安当}。他对她说：\Quote{我自己也想去，若不是被囚在这隐修院里，且我的前去能有所裨益的话。因为自前天起，人类已失去了那位真正是众人之父的人。}她相信了他的话，留在原地；几天后，传来圣安当已安眠的消息。\par

30. 或许有人会惊奇他所行的奇迹，或他那令人难以置信的斋戒、知识及谦逊。但最令我惊叹的，是他践踏尊荣与名声的能力。主教、司铎、成群的神职人员和隐修士，甚至基督徒主妇（这是一个极大的诱惑），以及来自城乡各处的乌合之众，都聚集在他周围，为要从他手中领受他所祝福的饼或油。但他心中只念着独处，以致有一天他决意离开，便弄来一头驴（他因斋戒几乎精疲力竭，几乎走不动），试图悄悄溜走。消息传遍四方，仿佛巴勒斯坦的荒凉被定为公共哀悼之日，成千上万不同年龄和性别的人聚拢来阻止他离开。他不为恳求所动，用棍杖敲打沙子，不断说：\Quote{我不愿使我的主成为欺骗者；我不能目睹教堂倾覆，基督的祭台被践踏，我儿子的血被倾倒。}所有在场的人开始明白，有一些隐秘的事启示给了他，他不愿承认，但他们依然看守着他，不让他走。于是，他决定并公开呼请大家作证：除非释放他，否则他不吃不喝。七天后，他才结束了斋戒；他告别了众多朋友，由无数群众陪着来到\Place{贝提利乌姆}。在那里，他说服群众回去，并挑选了四十位隐修士作为同伴，他们具有旅途所需的资源且能在斋戒期间（即日落之后）赶路。然后，他拜访了邻近旷野中住在名为\Place{吕克诺斯}地方的弟兄们，三天后来到\Place{特乌巴图斯}城堡，看望被流放在那里的宣信者主教\Person{德拉孔提乌斯}。主教因这样一位杰出人物的到来而无比欢欣。又过了三天，他出发前往\Place{巴比伦}，经过艰苦跋涉后到达。随后，他拜访了同样是宣信者的主教\Person{斐罗}；因为信奉\OriginalTerm{亚略异端}{Arian heresy}的皇帝君士坦提乌斯将两人都流放到了那些地区。离开那里后，他三天后到达了\Place{阿佛洛狄同}镇。在那里，他遇到了一位名叫\Person{拜萨内斯}的执事，此人饲养单峰驼，由于旷野中缺水，这些驼被雇来搭载那些想拜访圣安当的旅行者。他然后向弟兄们宣告，蒙福的圣安当的逝世周年纪念日即将来临，他必须在圣人去世的那个地方彻夜守夜。于是，经过三天穿越荒凉可怕的旷野，他们终于来到一座极高的山，在那里找到两位隐修士：\Person{依撒格}和\Person{培路西安努斯}，前者曾是圣安当的侍从之一。\par

31. 既然我们亲临其地，此时似乎是描述这位伟人住所的合宜时机。有一座高耸多石的山，绵延约一英里，山麓有泉涌出，泉水部分被沙吸收，部分流向平原，渐成溪流，两岸无数棕榈树荫蔽，增添了许多愉悦与魅力。在此地，人们常可看见这位长者与圣安当的门徒一同踱步。他们说，圣安当本人也曾在此歌唱、祈祷、劳作，并在疲倦时休息。那些藤蔓和灌木是他亲手所植；那座园圃是他亲自设计。这个浇灌园子的水池是他费尽辛劳挖成的。那把锄头他使用了多年。依拉良会躺在圣人的床上，仿佛余温尚存，深情地亲吻它。那间小室是方形的，每边仅有一人睡卧之长。此外，在高耸的山顶上，需沿蜿蜒小径艰难攀登而上，可见到两间同样大小的斗室，安当为避开成群访客或门徒的陪伴时，便居住其中。这些斗室是从岩石中凿出来的，仅装有门。当他们来到园中，依撒格说：\Quote{你看，这座园子有灌木和蔬菜；大约三年前，被一群野驴糟蹋了。安当抓住其中一头领头的驴，让它站住，用棍子抽打它的肋旁，问它为何吞吃自己没有播种的东西。从此以后，除了它们惯常来喝的水之外，它们再没有碰过任何东西，连一根灌木或蔬菜都没有。} 长者又请求带他去看安当的葬地，他们便把他领到一旁；但究竟是否给他看了坟墓，不得而知。据说，秘密行事是为了遵从安当的命令，并防止当地一位极其富有的人——佩尔加米乌斯——将圣人的遗体移到他家中，并为他建造一座纪念堂。\par

32. 他回到阿弗洛狄通后，只留两位弟兄在身边，在附近的旷野中居住，实行如此严格的斋戒与静默，以致他觉得自己此时才初次开始事奉基督。自从天闭不雨、大地遭受干旱以来，已过去三年；人们普遍说，连元素也在哀悼圣安当的去世。依拉良并没有比在其他地方更不为当地居民所知；面容憔悴、因饥饿而消瘦的男女们，恳切地哀求这位基督的仆人——作为圣安当的继承者——赐给他们雨水。依拉良看见他们，心中奇妙地生出怜悯，举目向天，举起双手，立刻获得了他们的请求。但说来奇怪，那片干涸多沙的地区，在降雨之后，竟意想不到地产生了大量蛇和毒物，以致许多被咬伤的人若不赶快跑到依拉良那里，就会立刻死去。于是，他祝福了些油，所有农夫和牧人用这油涂抹伤口，便获得了确实的医治。\par

33. 依拉良见即使在那个地方也有人向他表示惊人的敬意，便去了亚历山大，打算从那里渡海到更远的沙漠绿洲。因为他自度隐修生活以来从未在城中停留过，便转向布鲁基翁——距亚历山大不远，那里有他认识的几位弟兄，他们极其欣喜地接待了这位长者。到了夜里，他们忽然听见他的门徒在备驴，准备起程。于是他们俯伏在他脚下，恳求他不要离开；他们伏在门前，宣称宁愿死也不愿失去这样一位客人。他回答说：\Quote{我匆忙离去的原因，是怕连累你们。你们日后必定会发现，我不是无缘无故突然离开的。} 第二天，加沙的官员与行政长官的侍卫听说他前一天已到达，便进入隐修院；到处找不到他，便开始彼此说：\Quote{我们所听说的确是真的。他是个术士，能预知未来。} 事实上，在尤利安登基之后，依拉良离开巴勒斯坦、他的隐修院被毁之后，加沙城曾向皇帝呈上请愿书，要求处死依拉良和赫西基乌斯，并且到处张贴公告，命人搜捕他们。\par

34. 他离开布鲁基翁后，穿越无路的旷野进入绿洲，在那里住了大约一年。但既然他的名声也已传到那里，他感到自己在东方无法隐藏——许多人都因传闻和目睹认识他——于是开始考虑乘船去某个孤岛，好让自己既已暴露于陆地，至少能在海上找到藏身之处。正在那时，他的门徒哈德良从巴勒斯坦赶来，带来了尤利安已被弑杀、一位基督徒皇帝已开始统治的消息；因此，据说他应当回到他隐修院的遗骸那里。但他听到这话，庄严地拒绝了回去；他雇了一匹骆驼，穿越荒漠，到达利比亚海岸的城市帕雷托尼翁。在那里，不幸的哈德良想要返回巴勒斯坦，又不愿舍弃他的主人名声长久附带的荣光，便对他百般责骂，最后收拾好他带给他的主人的弟兄们的礼物，在依拉良不知情的情况下出发了。由于我今后没有机会再提到此人，我只愿记录——为警戒那些轻慢自己主人的人——不久之后，他患上了瘰疬，全身溃烂化为腐肉。\par

35. 老人由\Person{加匝诺}陪同，登上一艘驶往西西里的船。行至亚得里亚海中途，他正打算出售一部自己年轻时亲手抄写的\WorkTitle{福音书}来支付船费，船主的儿子被魔鬼附身，开始喊叫说：\Quote{天主之仆\Person{希拉律翁}，为什么借着你，我们在海上也不得安宁？暂且宽容我，等我上了岸。不要让我在这里被赶出去，抛入深渊。}圣人回答说：\Quote{如果我的天主允许你留下，你就留下；但如果祂把你赶走，为什么把憎恶加在我这罪人和乞丐身上呢？}他说这话，是为了使船上的水手和商人到了岸上不至于出卖他。不久之后，那孩子被洁净了，他的父亲和在场其他人都承诺决不向任何人透露圣人的名字。\par

36. 临近西西里的海岬\Place{帕基努斯}时，他把\WorkTitle{福音书}献给船主，作为自己和\Person{加匝诺}的船费。那人却不愿接受，尤其因为他看到除了那卷书和他们身上的衣服，他们一无所有；最后他发誓决不收取。但年迈的圣人，热忱而自信于自己的贫穷意识，因自己没有世俗财物，且被当地人视为乞丐而极其喜乐。\par

37. 他再次思考后，担心来自东方的商人会认出他，便逃往内陆，离海约二十英里。在那里一块废弃的地上，他每天捆一捆柴火，放在他门徒的背上，到附近的大宅出售。他们就以此维持生计，并能买一点面包给偶然来访的人。但正应验了经上所记的：\BibleRef{玛 5:14} \BibleQuote{建在山上的城是不能隐藏的。}恰巧有一个被魔鬼折磨的侍卫在罗马的圣伯多禄大殿里，那污鬼在他里面喊叫说：\Quote{几天前，基督之仆\Person{希拉律翁}进了西西里，没有人认识他，他以为自己是隐藏的。我要去出卖他。}他立即与随从登上一艘停泊在港口的船，航往\Place{帕基努斯}，被魔鬼引到老人的小屋，在那里俯伏在地，当场被治愈。这是他在西西里的第一个奇迹，使无数病人来到他那里（也带来了许多虔诚的人）；以至于一个患水肿肿胀的领袖，来到的当天就被治愈。后来他给圣人献上无数的礼物，但圣人用救主对门徒的话回答他：\BibleRef{玛 10:8} \BibleQuote{你们白白得来的，也要白白地施舍。}\par

38. 当此事在西西里发生时，他的门徒\Person{赫西基乌斯}正在世界各地寻找老人，遍历海岸，深入旷野，始终坚信无论他在哪里，都不可能长久隐藏。三年后，他在\Place{墨托涅}从一个卖旧衣服的犹太人那里听说，一位基督教的先知出现在西西里，大行奇迹异事，简直使人以为他是古代圣人之一。于是他询问那人的衣着、步态、言语，特别是年龄，但一无所知。他的消息提供者只是说他曾听说过此人。于是他渡过亚得里亚海，一路顺风到了\Place{帕基努斯}，在海湾岸边的一间小屋住下，打听老人的消息，从每个人讲述的故事中发现他在哪里、在做什么。最令众人惊讶的是，在行了如此伟大的奇迹异事后，他甚至没有从该地区的任何人那里接受过一块面包。长话短说，圣人\Person{赫西基乌斯}跪在师父膝前，用眼泪沾湿他的脚；最后师父温和地扶他起来，两三天谈论之后，他从\Person{加匝诺}得知，\Person{希拉律翁}觉得自己不能再住在那地方，而想要去某些野蛮部落，那里他的名号还不为人知。\par

39. 于是他带他到了达尔马提亚的\Place{埃皮达乌斯}城，他在附近乡间住了几天，但仍不能隐藏。一条巨蛇——当地人称之为\OriginalTerm{蟒蛇}{boas}，因为它们极大，常吞食牛——正在广袤的整个省份肆虐，不仅吞噬畜群，还吞噬那些被它呼吸之力吸过去的农夫和牧人。他命令为它准备一个火堆，然后向基督祈祷，召唤那爬虫，命它爬上柴堆，然后点燃了火。于是当众将那凶猛的野兽烧成灰烬。但他开始焦虑地问自己该怎么办，该往何处去。他再次准备逃跑，在思想中遍历荒凉之地，悲伤于他的奇迹能述说他，而他的舌头却沉默。\par

40. 当时，尤利安死后，全世界发生地震，致使海水漫过堤岸，船只被挂在悬崖边上。好似天主在恐吓第二次洪水，又像万物都要回归最初的洪荒。当\Place{埃庇道鲁斯}的居民看见这景象——咆哮的波涛、翻腾的海水、山一般高的漩涡冲击着海岸——他们担心其他地方发生的事也会降临到他们头上，使他们的城镇完全毁灭，便来到老人跟前，仿佛准备作战一般，把他安置在海岸上。他在沙上划了三次十字圣号，然后面向大海，伸出双手。谁也不会相信，那汹涌的海水竟如墙壁一般矗立在他面前。海水咆哮良久，仿佛对那屏障愤愤不平，然后渐渐消退到原来的水平线。\Place{埃庇道鲁斯}及周围地区至今传颂此事，母亲们教导子女，要子孙后代铭记。的确，那曾对宗徒们说的话：\Quote{你们若有信德，就能对这座山说：移到海里去，它必会移去。}这话甚至可以字面地应验，只要人有主命令宗徒们所拥有的那样的信德。因为，山沉入海中，与老人脚前那原本到处流淌的、山一般的巨大水墙突然变得像岩石一样坚硬，又有什么分别呢？\par

41. 全国都惊叹不已，这伟大奇迹的名声传遍各处，甚至传到了\Place{萨洛纳}。老人知道此事后，在夜间乘一艘小帆船秘密逃走；两天后找到一艘商船，来到\Place{塞浦路斯}。在\Place{马勒亚}和\Place{基西拉}之间，那些把桨帆并用舰队留在岸上的海盗，驾着两艘相当大的轻船向他们扑来；此外，他们四面八方还受到海浪的冲击。所有船夫都开始惊慌、哭泣、离开岗位、取出船桨，仿佛一条消息还不够，一再告诉老人海盗就在眼前。他远远地看了他们一眼，微微一笑，然后转向他的门徒说：\BibleRef{玛 14:32}\Quote{你们这小信德的人，为什么疑惑？这些人难道比法郎的军队还多？但他们全都因天主的旨意淹死了。}他这样说，但敌人并不退缩，船首冒着白沫，越来越近，现在只有一箭之遥。他站在船头，面向他们，伸出手说：\Quote{到此为止，不能再向前了。}奇妙的是，那两艘船立刻向后退，尽管被船桨推着向前，却越落越远。海盗们惊讶地发现自己正在后退，用尽全力想追上商船，但被带向海岸的速度比他们追来时快得多。\par

42. 其余的事情我略过不提，恐怕在我的历史中似乎在出版一卷奇迹。我只说这一点：当他在顺风中航行于\Place{基克拉泽斯}群岛之间时，他听到来自各处城镇乡村的邪魔喊叫，成群地跑到岸边。他随后进入\Place{塞浦路斯}的\Place{帕福斯}城，此城因诗人的诗歌而闻名，其神庙在多次地震后只剩下废墟，成为它昔日辉煌的唯一明证。他开始在离城约两英里处隐居，为能有几天的休息而高兴。但没过二十天，全岛凡有邪魔附身的人就开始呼喊：基督的仆人\Person{希拉良}来了，必须赶紧到他那里去。\Place{萨拉米斯}、\Place{库里翁}、\Place{拉佩塔}以及其他城市都同声呼喊；许多人声称他们认识\Person{希拉良}，他确实是基督的仆人，但不知他在哪里。于是，在短短三十天多一点的时间里，大约有两百人，男女都有，聚集到他那里。他看到他们时，哀叹他们不许他安静；他仿佛怀着某种复仇的渴望，用如此迫切的祈祷鞭策他们，以致一些人立即痊愈，另一些人在两三天后，所有人都在一周内得到了治愈。\par

43. 他在此住了两年，时常想着逃走。同时，他派赫西基乌斯在春天返回巴勒斯坦，去问候弟兄们并探访他的隐修院旧址。赫西基乌斯回来后，发现希拉里翁渴望再次乘船前往埃及，即那个叫布科利亚的地方；但赫西基乌斯劝说他，既然那里没有基督徒，只有凶猛的蛮族，倒不如去塞浦路斯本岛上一个更高更僻静的地方。经过长时间仔细搜寻，他找到这样一个地方：离海十二英里，位于崎岖山峦的深处，攀登时几乎要手足并用才能上去。他带领希拉里翁来到那里。老人进去环顾四周。这确实是一个荒凉可怕的地方：虽然四面树木环绕，山巅流水潺潺，一片宜人的小菜园，果树繁多（不过他从未吃过这些果子），但附近有一座古庙的废墟，据他自己常讲、弟子们也作证，从那里日夜传来无数魔鬼的回声，简直像有一支军队。他非常高兴想到仇敌就在附近，于是住了五年，在这最后的日子里，赫西基乌斯频繁的探访使他倍感欣慰；因为山路陡峭崎岖，加上无数鬼怪（传说如此），很少有人能或敢于上去见他。有一天，当他正要离开菜园时，看见一个完全瘫痪的人躺在门前。他问赫西基乌斯这人是谁，是怎么被带来的。赫西基乌斯回答说，那是这所菜园所属田庄的管家。希拉里翁痛哭不已，向躺在地上的人伸出手，说：\Quote{我奉我们主耶稣基督的名，命令你起来行走。}话还在他嘴边，那人四肢奇迹般地健壮起来，起来站得稳稳的。此事传开后，许多人的需要甚至战胜了无路的旅途和地方的险恶。周围所有房屋的住客最关心的事就是防止他逃走，因为有传言说他不会在一个地方久留。他的这种习惯并非出于轻浮或幼稚，而是因为他躲避众人的关注和赞誉，始终渴望静默和隐姓埋名的生活。\par

44. 在他八十岁那年，趁赫西基乌斯不在，他亲手写了一封短信作为遗嘱，把他所有的财富（即一本福音书、他的粗毛布长袍、风帽和斗篷）留给了赫西基乌斯，因为他的仆人几天前去世了。因此，许多热心人从帕福斯来到病人这里，特别是因为他们听说他说自己不久就要迁移到主那里，摆脱身体的束缚。还有一位圣妇君士坦提娅也来了，她的女婿和女儿曾被他用油敷抹而免于死亡。他恳切地请求他们所有人，死后不要让他停留片刻，而是立刻把他埋在同一个菜园里，就穿着他的山羊毛长袍、风帽和农夫斗篷。\par

45. 他的身体几乎冰冷，生命仅存理智。但他睁大眼睛反复说：\Quote{出去吧，你怕什么？出去吧，我的灵魂，你为何犹豫？你服侍基督将近七十年，还怕死亡吗？}说完这话，他便断了气。在城里人听到他死讯之前，他就被立即埋葬了。\par

46. 当圣人赫西基乌斯听到他去世的消息，便前往塞浦路斯。为了麻痹严加看守的当地人，他假装想在同一个菜园里居住；然后在大约十个月内，冒着生命危险，偷走了圣人的遗体。他把遗体带到马朱马；在那里，全体隐修士和成群的市民列队护送，将其安葬在古老的隐修院中。他的长袍、风帽和斗篷完好无损；整个身体如同活人一般完好，并散发着芬芳的香气，仿佛经过防腐处理。\par

47. 在我结束本书之际，我认为不应忽略提及圣妇君士坦提娅的虔诚。当她接到消息说希拉里翁的遗体在巴勒斯坦时，立即去世，甚至以死亡证明了她对天主仆人的真诚爱戴。因为她惯于整夜在他的墓前守夜，与他交谈，仿佛他就在眼前，以激励她的祈祷。直到今天，人们还可以看到巴勒斯坦人和塞浦路斯人之间奇特的争论：一方声称拥有希拉里翁的遗体，另一方则声称拥有他的精神。然而，两地每天都有许多奇迹发生，但在塞浦路斯的菜园里更多，或许因为那块地方是他最钟爱的。\par

\SourceRecord{Translated by W.H. Fremantle, G. Lewis and W.G. Martley.；Nicene and Post-Nicene Fathers, Second Series；Vol. 6.；Edited by Philip Schaff and Henry Wace.；Buffalo, NY: Christian Literature Publishing Co.,；1893.；<http://www.newadvent.org/fathers/3003.htm>.}
